\chapter{高级类 Schnorr 签名 (Advanced Schnorr-like Signatures)}
\label{chapter:advanced-schnorr}

一个基本的 Schnorr 签名只有一个签名密钥。事实上，我们可以应用其核心概念来创建各种复杂度递增的签名方案。其中一种方案 MLSAG 将在 Monero 的 transaction 协议中具有核心重要性。



%-same signing key across multiple bases
\section{跨多个基底证明离散对数知识 (Prove knowledge of a discrete logarithm across multiple bases)}
\label{sec:proofs-discrete-logarithm-multiple-bases}

证明同一个私钥被用于在不同"基底"密钥上构建公钥通常很有用。例如，我们可以有一个普通公钥 $k G$，以及与某个人的公钥的 Diffie-Hellman 共享密钥 $k R$（回顾第 \ref{DH_exchange_section} 节），其中基底密钥是 $G$ 和 $R$。正如我们将很快看到的，我们可以证明知道 $k G$ 中的离散对数 $k$，证明知道 $k R$ 中的 $k$，{\em 并}证明两种情况下的 $k$ 是相同的（全部不泄露 $k$）。


\subsection*{非交互式证明 (Non-interactive proof)}

假设我们有一个私钥 $k$ 和 $d$ 个基底密钥 $\mathcal{J} = \{J_1,...,J_d\}$。对应的公钥为 $\mathcal{K} = \{K_1,...,K_d\}$。我们在所有基底上进行类 Schnorr 证明（回顾第 \ref{sec:schnorr-fiat-shamir} 节）。\footnote{虽然我们说的是"证明"，但通过在挑战 hash 中包含消息 $\mathfrak{m}$ 可以轻易地将其变为签名。在此语境中，这两个术语可以互换使用。}假设存在一个 hash 函数 \(\mathcal{H}_n\)
映射到0到 $l-1$ 之间的整数。\footnote{在 Monero 中，hash 函数 $\mathcal{H}_n(x) = \textrm{sc\textunderscore reduce32}(\mathit{Keccak}(x))$，其中 $\mathit{Keccak}$ 是 SHA3 的基础，sc\textunderscore reduce32() 将256位结果放在0到 $l-1$ 的范围内（虽然严格来说应该是1到 $l-1$）。}

\begin{enumerate}
	\item 生成随机数 $\alpha \in_R \mathbb{Z}_l$，并计算所有 $i \in (1,...,d)$ 的 $\alpha J_i$。
	\item 计算挑战，\vspace{.175cm}
	\[c = \mathcal{H}_n(\mathcal{J},\mathcal{K},[\alpha J_1],[\alpha J_2],...,[\alpha J_d])\]
	\item 定义响应 $r = \alpha - c*k$。
	\item 公开签名 $(c, r)$。
\end{enumerate}


\subsection*{验证 (Verification)}

假设验证者知道 $\mathcal{J}$ 和 $\mathcal{K}$，他执行以下操作。

\begin{enumerate}
	\item 计算挑战：\vspace{.175cm}
	\[c' = \mathcal{H}(\mathcal{J},\mathcal{K},[r J_1 + c*K_1],[r J_2 + c*K_2],...,[r J_d + c*K_d])\]
	\item 如果 $c = c'$，则签名者必须知道跨所有基底的离散对数，且每种情况下都是同一个离散对数（一如既往，除了可忽略的概率外）。
\end{enumerate}


\subsection*{为什么有效 (Why it works)}

如果 $d$ 个基底密钥只有一个，这个证明显然与我们原始的 Schnorr 证明相同（第 \ref{sec:schnorr-fiat-shamir} 节）。我们可以将每个基底密钥单独来看，发现多基底证明只是一系列连接在一起的 Schnorr 证明。此外，通过对所有这些证明使用唯一的挑战和响应，它们必须具有相同的离散对数 $k$。要获得适用于多个密钥的单一响应，挑战需要在为每个密钥定义 $\alpha$ 之前就被知道，但 $c$ 是 $\alpha$ 的函数！



%-multiple signing keys on their own unique bases (or they can be the same e.g. G)
\section{一个证明中的多个私钥 (Multiple private keys in one proof)}
\label{sec:multiple_private_keys_in_one_proof}

与多基底证明非常相似，我们可以组合多个使用不同私钥的 Schnorr 证明。这样做证明了我们知道一组公钥的所有私钥，并通过为所有证明只生成一个挑战来减少存储需求。


\subsection*{非交互式证明 (Non-interactive proof)}

假设我们有 $d$ 个私钥 $k_1,...,k_d$，以及基底密钥 $\mathcal{J} = \{J_1,...,J_d\}$。\footnote{这里没有理由 $\mathcal{J}$ 不能包含重复的基底密钥，或者所有基底密钥都相同（例如 $G$）。对于多基底证明来说重复是冗余的，但现在我们处理的是不同的私钥。}对应的公钥为 $\mathcal{K} = \{K_1,...,K_d\}$。我们同时为所有密钥进行类 Schnorr 证明。

\begin{enumerate}
	\item 为所有 $i \in (1,...,d)$ 生成随机数 $\alpha_i \in_R \mathbb{Z}_l$，并计算所有 $\alpha_i J_i$。
	\item 计算挑战，\vspace{.175cm}
	\[c = \mathcal{H}_n(\mathcal{J},\mathcal{K},[\alpha_1 J_1],[\alpha_2 J_2],...,[\alpha_d J_d])\]
	\item 定义每个响应 $r_i = \alpha_i - c*k_i$。
	\item 公开签名 $(c, r_1,...,r_d)$。
\end{enumerate}


\subsection*{验证 (Verification)}

假设验证者知道 $\mathcal{J}$ 和 $\mathcal{K}$，他执行以下操作。

\begin{enumerate}
	\item 计算挑战：\vspace{.175cm}
	\[c' = \mathcal{H}(\mathcal{J},\mathcal{K},[r_1 J_1 + c*K_1],[r_2 J_2 + c*K_2],...,[r_d J_d + c*K_d])\]
	\item 如果 $c = c'$，则签名者必须知道 $\mathcal{K}$ 中所有公钥的私钥（除了可忽略的概率外）。
\end{enumerate}



\section{自发匿名群 (SAG) 签名 (Spontaneous Anonymous Group (SAG) signatures)}
\label{SAG_section}

群签名是一种证明签名者属于某个群组而不必识别他的方式。最初（Chaum 在 \cite{Chaum:1991:GS:1754868.1754897} 中），群签名方案要求系统由可信人员设置，并在某些情况下由其管理，以防止非法签名，并在少数方案中裁决争议。这些方案依赖于{\em 群密钥}，这是不理想的，因为它造成了可能破坏匿名性的泄露风险。此外，要求群成员之间的协调（即设置和管理）对于小群组或公司内部以外的场景是不可扩展的。

Liu {\em 等人}在 \cite{Liu2004} 中基于 Rivest {\em 等人} 在 \cite{rivest-leak-secret} 中的工作提出了一个更有趣的方案。作者详细介绍了一种名为 LSAG 的群签名算法，具有三个特性：{\em 匿名性、可链接性}和{\em 自发性}。这里我们讨论 SAG，即 LSAG 的不可链接版本，以获得概念上的清晰性。我们将可链接性的概念留到后面的章节。

具有匿名性和自发性的方案通常被称为"环签名"。在 Monero 的语境中，它们最终将实现不可伪造的、签名者模糊的 transaction，使资金流向在很大程度上不可追踪。


\subsection*{签名 (Signature)}

环签名由一个环和一个签名组成。每个{\em 环}是一组公钥，其中一个属于签名者，其余的与之无关。{\em 签名}使用该密钥环生成，任何验证它的人都无法分辨哪个环成员是实际的签名者。

我们在第 \ref{sec:signing-messages} 节中的类 Schnorr 签名方案可以被视为单密钥环签名。通过不立即定义 $r$，而是生成一个诱饵 $r'$ 并创建新的挑战来定义 $r$，我们得到了双密钥版本。
\\

设 \(\mathfrak{m}\) 为待签名的消息，\(\mathcal{R} = \{K_1, K_2, ..., K_n\}\) 为一组不重复的公钥（一个群/环），\(k_\pi\) 为签名者的私钥，对应其公钥 \(K_\pi \in \mathcal{R}\)，其中 $\pi$ 是秘密索引。

\begin{enumerate}
	\item 生成随机数 \(\alpha \in_R \mathbb{Z}_l\) 和伪响应 \(r_i \in_R \mathbb{Z}_l\)，其中 \(i \in \{1, 2, ..., n\}\) 但排除 \(i = \pi\)。

	\item 计算
	\[c_{\pi+1} = \mathcal{H}_n(\mathcal{R}, \mathfrak{m}, [\alpha G])\]

	\item 对于 \(i = \pi+1, \pi+2, ..., n, 1, 2, ..., \pi-1\) 计算，将 \(n + 1 \rightarrow 1\)，\vspace{.175cm}
	\[  c_{i+1} = \mathcal{H}_n(\mathcal{R}, \mathfrak{m}, [r_i G + c_i K_i])\] 

	\item 定义真实响应 $r_\pi$ 使得 \(\alpha = r_\pi + c_\pi k_\pi \pmod l\)。
\end{enumerate}

环签名包含签名 \(\sigma(\mathfrak{m}) = (c_1, r_1, ..., r_n) \) 和环 $\mathcal{R}$。


\subsection*{验证 (Verification)}

验证意味着证明 $\sigma(\mathfrak{m})$ 是由 $\mathcal{R}$ 中某个公钥对应的私钥创建的有效签名（不必知道是哪一个），方式如下：

\begin{enumerate}
	\item 对于 \(i = 1, 2, ..., n\) 迭代计算，将 \(n + 1 \rightarrow 1\)，\vspace{.175cm}
	\begin{align*}
	c'_{i+1}   = \mathcal{H}_n(\mathcal{R}, \mathfrak{m}, [r_i G + c_i {K_i}])
	\end{align*}

	\item 如果 \(c_1 = c'_1\)，则签名有效。注意 $c'_1$ 是最后计算的项。
\end{enumerate}

在此方案中，我们存储 (1+$n$) 个整数并使用 $n$ 个公钥。


\subsection*{为什么有效 (Why it works)}

我们可以通过一个例子来非正式地说服自己该算法是有效的。考虑环 $R = \{K_1, K_2, K_3\}$，其中 $k_\pi = k_2$。首先签名：
\begin{enumerate}
    \item 生成随机数：$\alpha$，$r_1$，$r_3$
\begin{align*}
    \intertext{\item 初始化签名循环：}	c_3 &= \mathcal{H}_n(\mathcal{R}, \mathfrak{m}, [\alpha G])
    \intertext{\item 迭代：\vspace{-.2cm}}
        c_1 &= \mathcal{H}_n(\mathcal{R}, \mathfrak{m}, [r_3 G + c_3 K_3])\\
        c_2 &= \mathcal{H}_n(\mathcal{R}, \mathfrak{m}, [r_1 G + c_1 K_1])
\end{align*}
    \item 通过响应来闭合循环：$r_2 = \alpha - c_2 k_2 \pmod{l}$
\end{enumerate}

我们可以将 $\alpha$ 代入 $c_3$ 来看看"环"这个词的由来：\vspace{.175cm}
\begin{alignat*}{3}
    c_3 &= \mathcal{H}_n(\mathcal{R}, \mathfrak{m}, [(r_2 + c_2 k_2) G &&])\\
    c_3 &= \mathcal{H}_n(\mathcal{R}, \mathfrak{m}, [r_2 G + c_2 K_2 &&])
\end{alignat*}\vspace{.05cm}

然后使用 $\mathcal{R}$ 和 $\sigma(\mathfrak{m}) = (c_1, r_1, r_2, r_3)$ 进行验证：
\begin{enumerate}
    \item 我们使用 $r_1$ 和 $c_1$ 来计算\vspace{.175cm}
    \begin{align*}
c'_2 &= \mathcal{H}_n(\mathcal{R}, \mathfrak{m}, [r_1 G + c_1 K_1])
    \intertext{\item 根据创建签名时的情况，我们看到 $c'_2 = c_2$。使用 $r_2$ 和 $c'_2$ 我们计算\vspace{.175cm}}
c'_3 &= \mathcal{H}_n(\mathcal{R}, \mathfrak{m}, [r_2 G + c'_2 K_2])
    \intertext{\item 通过将 $c_2$ 替换为 $c'_2$，我们可以很容易看出 $c'_3 = c_3$。使用 $r_3$ 和 $c'_3$ 我们得到\vspace{.175cm}}
c'_1 &= \mathcal{H}_n(\mathcal{R}, \mathfrak{m}, [r_3 G + c'_3 K_3])
    \end{align*}
\end{enumerate}
\quad 不出意外：如果我们将 $c_3$ 替换为 $c'_3$，则 $c'_1 = c_1$。\vspace{-.3cm}



\section{Back 的可链接自发匿名群 (bLSAG) 签名 (Back's Linkable Spontaneous Anonymous Group (bLSAG) signatures)}
\label{blsag_note}

从这里开始讨论的环签名方案展示了几个对产生 confidential transaction 有用的特性。\footnote{请记住，所有健壮的签名方案都有包含各种特性的安全模型。这里提到的特性也许与理解 Monero 环签名的目的最为相关，但不是对可链接环签名特性的全面概述。}注意，"签名者模糊性"和"不可伪造性"也适用于 SAG 签名。

\begin{description}
	\item[签名者模糊性 (Signer Ambiguity)]
	观察者应该能够确定签名者必须是环的成员（除了可忽略的概率外），但不能确定是哪个成员。\footnote{\label{anonymity_note}操作的匿名性通常用"匿名集"来衡量，即"所有可能执行该操作的人"。最大的匿名集是"全人类"，对于 Monero 来说是环的大小，例如所谓的"混合级别" $v$ 加上真实签名者。混合（Mixin）指的是每个环签名有多少个假成员。如果混合为 $v$ = 4，则有5个可能的签名者。扩大匿名集使得追踪真实行为者变得越来越困难。}Monero 用此来混淆每笔 transaction 中资金的来源。

	\item[可链接性 (Linkability)]
	如果一个私钥被用来签署两个不同的消息，那么这些消息将被链接在一起。\footnote{\label{linkability_note}可链接性属性不适用于非签名公钥。也就是说，公钥已被用于不同环签名的环成员不会导致链接。}正如我们将要展示的，此属性用于防止 Monero 中的双花攻击（除了可忽略的概率外）。

	\item[不可伪造性 (Unforgeability)]
    没有攻击者能够在除了可忽略的概率外伪造签名。\footnote{\label{unforgeability_note}某些环签名方案，包括 Monero 中的方案，对于自适应选择消息和自适应选择公钥攻击具有强安全性。攻击者即使能够获取所选消息的合法签名，并且这些签名对应于他选择的环中的特定公钥，也无法发现如何伪造哪怕一条消息的签名。这被称为{\em 存在不可伪造性}；参见 \cite{MRL-0005-ringct} 和 \cite{Liu2004}。}这用于防止没有相应私钥的人盗窃 Monero 资金。
\end{description}

在 LSAG 签名方案 \cite{Liu2004} 中，私钥的所有者可以为每个环产生一个匿名的、不可链接的签名。\footnote{\label{lsag_linkability_note}在 LSAG 方案中，可链接性仅适用于使用相同成员且相同顺序的环的签名，即"完全相同的环"。实际上是"每个环成员每个环一个匿名签名"。即使为不同的消息制作，签名也可以被链接。}在本节中，我们提出了 LSAG 算法的增强版本，其中可链接性独立于环的诱饵成员。\footnote{LSAG 在本报告第一版中有讨论。\cite{ztm-1}}

该修改在 \cite{MRL-0005-ringct} 中被揭示，基于 Adam Back \cite{AdamBack-ring-efficiency} 关于 CryptoNote \cite{cryptoNoteWhitePaper} 环签名算法的出版物（此前在 Monero 中使用，现已弃用；参见第 \ref{subsec:proofs-input-creation-spendproof} 节），而这又受到 Fujisaki 和 Suzuki 在 \cite{Fujisaki2007} 中工作的启发。


\subsection*{签名 (Signature)}

与 SAG 一样，设 \(\mathfrak{m}\) 为待签名的消息，\(\mathcal{R} = \{K_1, K_2, ..., K_n\}\) 为一组不重复的公钥，\(k_\pi\) 为签名者的私钥，对应其公钥 \(K_\pi \in \mathcal{R}\)，其中 $\pi$ 是秘密索引。假设存在一个 hash 函数 \(\mathcal{H}_p\)，映射到 EC 中的曲线点。\footnote{$\mathcal{H}_p$ 产生的点是否压缩并不重要。它们总是可以被解压的。}\footnote{Monero 使用一个\marginnote{src/ringct/ rctOps.cpp {\tt hash\_to\_p3()}}直接返回曲线点的 hash 函数，而不是计算某个整数然后乘以 $G$。如果有人发现了找到 $n_x$ 使得 $n_x G = \mathcal{H}_p(x)$ 的方法，$\mathcal{H}_p$ 就会被破解。参见 \cite{hashtopoint-writeup} 中的算法描述。根据 CryptoNote 白皮书 \cite{cryptoNoteWhitePaper}，其来源是这篇论文：\cite{hashtopoint-original-paper}。}

\begin{enumerate}
	\item 计算 key image \(\tilde{K} = k_\pi \mathcal{H}_p(K_\pi)\)。\footnote{在 Monero 中，使用 hash 到点函数来计算 key image 而不是另一个基底点非常重要，这样线性性就不会导致由相同地址（即使是一次性地址）创建的签名被链接。参见 \cite{cryptoNoteWhitePaper} 第18页。}

	\item 生成随机数 \(\alpha \in_R \mathbb{Z}_l\) 和随机数 \(r_i \in_R \mathbb{Z}_l\)，其中 \(i \in \{1, 2, ..., n\}\) 但排除 \(i = \pi\)。

	\item 计算
	\[c_{\pi+1} = \mathcal{H}_n(\mathfrak{m}, [\alpha G], [\alpha \mathcal{H}_p(K_\pi)])\]

	\item 对于 \(i = \pi+1, \pi+2, ..., n, 1, 2, ..., \pi-1\) 计算，将 \(n + 1 \rightarrow 1\)，\vspace{.175cm}
	\[c_{i+1} = \mathcal{H}_n(\mathfrak{m}, [r_i G + c_i K_i], [r_i \mathcal{H}_p(K_i) + c_i \tilde{K}])\]

	\item 定义 \(r_\pi = \alpha - c_\pi k_\pi \pmod l\)。
\end{enumerate}

签名为 \(\sigma(\mathfrak{m}) = (c_1, r_1, ..., r_n)\)，附带 key image $\tilde{K}$ 和环 $\mathcal{R}$。


\subsection*{验证 (Verification)}

验证意味着证明 $\sigma(\mathfrak{m})$ 是由 $\mathcal{R}$ 中某个公钥对应的私钥创建的有效签名，方式如下：

\begin{enumerate}
    \item 检查 $l \tilde{K} \stackrel{?}{=} 0$。
	\item 对于 \(i = 1, 2, ..., n\) 迭代计算，将 \(n + 1 \rightarrow 1\)，\vspace{.175cm}
	\begin{align*}
	c'_{i+1} = \mathcal{H}_n(\mathfrak{m}, [r_i G + c_i {K_i}], [r_i \mathcal{H}_p(K_i) + c_i \tilde{K}])
	\end{align*}

	\item 如果 \(c_1 = c'_1\)，则签名有效。
\end{enumerate}

在此方案中，我们存储 (1+$n$) 个整数，有一个 EC key image，并使用 $n$ 个公钥。

我们\marginnote{src/crypto- note\_core/ cryptonote\_ core.cpp {\tt check\_tx\_ inputs\_key- images\_do- main()}}必须检查 $l \tilde{K} \stackrel{?}{=} 0$，因为可以将大小为 $h$（余因子）的子群中的 EC 点加到 $\tilde{K}$ 上，并且如果所有 $c_i$ 都是 $h$ 的倍数（我们可以使用不同的 $\alpha$ 和 $r_i$ 值通过自动试错来实现），就可以使用相同的环和签名密钥制作 $h$ 个不可链接的有效签名。\footnote{我们不关心来自其他子群的点，因为 $\mathcal{H}_n$ 的输出限制在 $\mathbb{Z}_l$ 中。对于 EC 阶 $N = h l$，$N$ 的所有因数（以及因此可能的子群）要么是 $l$（素数）的倍数，要么是 $h$ 的因数。}这是因为 EC 点乘以其子群的阶等于零。\footnote{在 Monero 的早期历史中，这没有被检查。幸运的是，在2017年4月修复实现（协议第5版）之前，这个漏洞没有被利用 \cite{key-image-bug}。}

明确地说，给定阶为 $l$ 的子群中的某个点 $K$，阶为 $h$ 的某个点 $K^h$，以及一个能被 $h$ 整除的整数 $c$：
\begin{align*}
    c*(K + K^h) &= cK + cK^h\\
                &= cK + 0
\end{align*}

我们可以用与更简单的 SAG 签名方案类似的方式展示正确性（即"它是如何工作的"）。

我们的描述力求忠实于 bLSAG 的原始解释，该解释在计算 $c_i$ 的 hash 中不包含 $\mathcal{R}$。将密钥包含在 hash 中被称为"密钥前缀"。最近的研究 \cite{key-prefix-paper} 表明这可能不是必需的，尽管添加前缀是类似签名方案的标准做法（LSAG 使用密钥前缀）。


\subsection*{可链接性 (Linkability)}

给定两个在某些方面不同的有效签名（例如不同的伪响应、不同的消息、不同的整体环成员），\vspace{.1cm}
\begin{align*}
	\sigma(\mathfrak{m})   &= (c_1, r_1, ..., r_n)\textrm{ with } \tilde{K}\textrm{, and}\\
	\sigma'(\mathfrak{m}')  &= (c_1', r'_1, ..., r'_{n'})\textrm{ with } \tilde{K}'\textrm{,}
\end{align*}
\quad 如果 \(\tilde{K} =  \tilde{K}'\)，则显然两个签名来自同一个私钥。

虽然观察者可以将 $\sigma$ 和 $\sigma'$ 链接起来，但他不一定知道 $\mathcal{R}$ 或 $\mathcal{R}'$ 中的哪个 $K_i$ 是罪魁祸首，除非它们之间只有一个共同密钥。如果有多个共同环成员，他唯一的办法就是求解 DLP 或以某种方式审计环（例如学习所有 $i \neq \pi$ 的 $k_i$，或学习 $k_\pi$）。\footnote{\label{lsag_unforgeable_note}LSAG 与 bLSAG 非常相似，是不可伪造的，这意味着没有攻击者可以在不知道私钥的情况下制作有效的环签名。如果他发明了一个假的 $\tilde{K}$ 并用 $c_{\pi+1}$ 初始化他的签名计算，那么，由于不知道 $k_\pi$，他无法计算一个数 $r_\pi = \alpha - c_\pi k_\pi$ 来产生 $[r_\pi G + c_\pi K_\pi] = \alpha G$。验证者会拒绝他的签名。Liu {\em 等人}证明，能够通过验证的伪造行为概率极低 \cite{Liu2004}。}



\section{多层可链接自发匿名群 (MLSAG) 签名 (Multilayer Linkable Spontaneous Anonymous Group (MLSAG) signatures)}
\label{sec:MLSAG}

为了签署 transaction，必须使用多个私钥进行签名。在 \cite{MRL-0005-ringct} 中，Shen Noether {\em 等人}描述了一种 bLSAG 签名方案的多层推广，适用于我们有 \(n \cdot m\) 个密钥集合的情况；即集合\vspace{.175cm}
\[\mathcal{R} = \{K_{i,j}\}  \quad \textrm{for} \quad  i \in \{1, 2, ..., n\} \quad \textrm{and} \quad j \in \{1, 2, ..., m\}\]

其中我们知道对应于子集 \(\{K_{\pi, j}\}\) 的 $m$ 个私钥 \(\{k_{\pi, j}\}\)，对于某个索引 \(i = \pi\)。如果我们推广可链接性的概念，这样的算法将满足我们的需求。
\begin{description}
	\item[可链接性 (Linkability)] 如果任何私钥 \(k_{\pi, j}\) 被用于2个不同的签名中，则这些签名将被自动链接。
\end{description}


\subsection*{签名 (Signature)}

\begin{enumerate}
	\item 为所有 \(j \in \{1, 2, ..., m\}\) 计算 key image \(\tilde{K_j} = k_{\pi, j} \mathcal{H}_p(K_{\pi, j})\)。

	\item 生成随机数 \(\alpha_j \in_R \mathbb{Z}_l\)，以及 \(r_{i, j} \in_R \mathbb{Z}_l\)，其中 \(i \in \{1, 2, ..., n\}\)（排除 \(i = \pi\)）和 \(j \in \{1, 2, ..., m\}\)。

	\item 计算\footnote{Monero\marginnote{src/ringct/ rctSigs.cpp {\tt MLSAG\_Gen()}} MLSAG 使用密钥前缀。每个挑战包含如下显式公钥（添加了 bLSAG 中缺失的 $K$ 项；key image 包含在签名的消息中）：\vspace{-.25cm}
	\[c_{\pi+1} = \mathcal{H}_n(\mathfrak{m}, K_{\pi, 1}, [\alpha_1 G], [\alpha_1 \mathcal{H}_p(K_{\pi, 1})], ..., K_{\pi, m}, [\alpha_m G], [\alpha_m \mathcal{H}_p(K_{\pi, m})])
	\]}
	\[c_{\pi+1} = \mathcal{H}_n(\mathfrak{m}, [\alpha_1 G], [\alpha_1 \mathcal{H}_p(K_{\pi, 1})], ..., [\alpha_m G], [\alpha_m \mathcal{H}_p(K_{\pi, m})])\]

	\item 对于 \(i = \pi+1, \pi+2, ..., n, 1, 2, ..., \pi-1\) 计算，将 \(n + 1 \rightarrow 1\)，\vspace{.175cm}
	\[ c_{i+1} = \mathcal{H}_n(\mathfrak{m}, [r_{i, 1} G + c_i K_{i, 1}], [r_{i, 1} \mathcal{H}_p(K_{i, 1}) + c_i \tilde{K}_1], 
	..., [r_{i, m} G + c_i K_{i, m}], [r_{i, m} \mathcal{H}_p(K_{i, m}) + c_i \tilde{K}_m])\]

	\item 定义所有 \(r_{\pi, j} = \alpha_j - c_\pi k_{\pi, j} \pmod l\)。
\end{enumerate}

签名为 \(\sigma(\mathfrak{m}) = (c_1, r_{1, 1}, ..., r_{1, m}, ..., r_{n, 1}, ..., r_{n, m}) \)，附带 key image $(\tilde{K}_1, ...,  \tilde{K}_m)$。

%理解 MLSAG 的一种方式是：有 $m$ 个大小为 $n$ 的子循环，在每个子循环中我们在索引 $i = \pi$ 处知道一个私钥（总共 $m \cdot n$ 个公钥）。签名算法在每个阶段 $c$ 将一个"栈"的密钥绑定在一起，由每个子循环中的一个密钥组成。bLSAG 是 $m = 1$ 的特殊情况。


\subsection*{验证 (Verification)}

签名的验证按以下方式进行：

\begin{enumerate}
    \item 对所有 $j \in \{1,...,m\}$ 检查 $l \tilde{K}_j \stackrel{?}{=} 0$。
	\item 对于 \(i = 1, ..., n\) 计算，将 \(n + 1 \rightarrow 1\)，\vspace{.175cm}
	\begin{align*}
	c'_{i+1} = \mathcal{H}_n(\mathfrak{m}, [r_{i, 1} G + c_i K_{i, 1}], [r_{i, 1} \mathcal{H}_p(K_{i, 1}) + c_i \tilde{K}_1], 
	..., [r_{i, m} G + c_i K_{i, m}], [r_{i, m} \mathcal{H}_p(K_{i, m}) + c_i \tilde{K}_m])
	\end{align*}

	\item 如果 \(c_1 = c'_1\)，则签名有效。
\end{enumerate}


\subsection*{为什么有效 (Why it works)}

正如 SAG 算法一样，我们可以容易地观察到

\begin{itemize}
    \item 如果 \(i \ne \pi \)，则显然值 \(c'_{i + 1}\) 是按照签名算法中描述的方式计算的。

    \item 如果 \(i = \pi\)，由于 \(r_{\pi, j} = \alpha_j - c_\pi k_{\pi, j} \) 闭合了循环，\vspace{.175cm}
    \begin{alignat*}{6}
        r_{\pi, j} G + c_\pi K_{\pi,j} &= (\alpha_j - c_\pi k_{\pi, j}) G + c_\pi K_{\pi,j} = \alpha_j G\\
        \intertext{并且}
        r_{\pi, j} \mathcal{H}_p(K_{\pi, j}) + c_\pi \tilde{K}_j &= (\alpha_j - c_\pi k_{\pi, j}) \mathcal{H}_p(K_{\pi, j}) + c_\pi \tilde{K}_j = \alpha_j \mathcal{H}_p(K_{\pi, j})\\
    \end{alignat*}
    换句话说，\(c'_{\pi + 1} = c_{\pi+1}\) 也成立。
\end{itemize}


\subsection*{可链接性 (Linkability)}

如果私钥 \(k_{\pi, j}\) 被重复使用来制作任何签名，签名中提供的对应 key image \(\tilde{K}_j\) 将揭示这一点。这一观察符合我们对可链接性的推广定义。\footnote{与 bLSAG 一样，链接的 MLSAG 签名不会指明使用了哪个公钥进行签名。然而，如果链接的 key image 的子循环环只有一个共同密钥，那么罪魁祸首就是显而易见的。如果罪魁祸首被识别，两个签名的所有其他签名成员都会被揭露，因为它们共享罪魁祸首的索引。}


\subsection*{空间需求 (Space requirements)}

在此方案中，我们存储 (1+$m*n$) 个整数，有 $m$ 个 EC key image，并使用 $m*n$ 个公钥。



\section{简洁可链接自发匿名群 (CLSAG) 签名 (Concise Linkable Spontaneous Anonymous Group (CLSAG) signatures)}
\label{sec:CLSAG}

CLSAG \cite{MRL-0011-CLSAG}\footnote{本节所依据的论文是一份正在定稿以供外部评审的预印本。CLSAG 作为 MLSAG 在未来协议版本中的替代方案很有前景，但尚未实现，也可能不会在未来实现。}是介于 bLSAG 和 MLSAG 之间的一种方案。假设你有一个"主"密钥，与之相关联的是几个"辅助"密钥。证明知道所有私钥很重要，但可链接性仅适用于主密钥。这种可链接性收缩允许比 MLSAG 更小、更快的签名。

与 MLSAG 一样，我们有一个 \(n \cdot m\) 个密钥的集合（$n$ 是环的大小，$m$ 是签名密钥的数量），主密钥位于索引1。换句话说，有 $n$ 个主密钥，第 $\pi$ 个主密钥及其辅助密钥将进行签名。\vspace{.175cm}
\[\mathcal{R} = \{K_{i,j}\}  \quad \textrm{for} \quad  i \in \{1, 2, ..., n\} \quad \textrm{and} \quad j \in \{1, 2, ..., m\}\]

我们知道对应于子集 \(\{K_{\pi, j}\}\) 的私钥 \(\{k_{\pi, j}\}\)，对于某个索引 \(i = \pi\)。


\subsection*{签名 (Signature)}

\begin{enumerate}
	\item 为所有 \(j \in \{1, 2, ..., m\}\) 计算 key image \(\tilde{K_j} = k_{\pi, j} \mathcal{H}_p(K_{\pi, 1})\)。注意基底密钥始终相同，因此 $j>1$ 的 key image 是"辅助 key image"。为简化符号，我们统称它们为 $\tilde{K}_j$。

	\item 生成随机数 \(\alpha \in_R \mathbb{Z}_l\)，以及 \(r_{i} \in_R \mathbb{Z}_l\)，其中 \(i \in \{1, 2, ..., n\}\)（排除 \(i = \pi\)）。

    \item 为 \(i \in \{1, 2, ..., n\}\) 计算聚合公钥 $W_i$ 和聚合 key image $\tilde{W}$\footnote{CLSAG 论文建议使用不同的 hash 函数来进行域分离，我们通过在每个 hash 前加标签来建模 \cite{MRL-0011-CLSAG}，例如 $T_1 =$ "CLSAG\_1"，$T_c =$ "CLSAG\_c" 等。域分离的 hash 函数即使输入相同也产生不同的输出。我们这里也使用密钥前缀（将包含所有密钥的 $\mathcal{R}$ 包含在 hash 中）。域分离是 Monero 开发的新策略，未来添加的所有 hash 函数应用（v13+）都可能会这样做。历史上的 hash 函数使用可能将保持不变。}%实际实现可能使用 $j$ 标志，而不是索引本身
    \begin{align*}
    W_i &= \sum^{m}_{j=1} \mathcal{H}_n(T_j, \mathcal{R}, \tilde{K}_1,...,\tilde{K}_{m})*K_{i,j}\\
    \tilde{W} &= \sum^{m}_{j=1} \mathcal{H}_n(T_j, \mathcal{R}, \tilde{K}_1,...,\tilde{K}_{m})*\tilde{K}_j
    \end{align*}{}
    其中 $w_{\pi} = \sum_j \mathcal{H}_n(T_j,...)*k_{\pi,j}$ 是聚合私钥。


	\item 计算
	\[c_{\pi+1} = \mathcal{H}_n(T_c, \mathcal{R}, \mathfrak{m}, [\alpha G], [\alpha \mathcal{H}_p(K_{\pi, 1})])\]

	\item 对于 \(i = \pi+1, \pi+2, ..., n, 1, 2, ..., \pi-1\) 计算，将 \(n + 1 \rightarrow 1\)，\vspace{.175cm}
	\[c_{i+1} = \mathcal{H}_n(T_c, \mathcal{R}, \mathfrak{m}, [r_i G + c_i W_i], [r_{i} \mathcal{H}_p(K_{i,1}) + c_i \tilde{W}])\]

	\item 定义 \(r_{\pi} = \alpha - c_\pi w_\pi \pmod l\)。
\end{enumerate}

因此 \(\sigma(\mathfrak{m}) = (c_1, r_1, ..., r_n) \)，附带主 key image $\tilde{K}_1$ 和辅助 image $(\tilde{K}_2,...,\tilde{K}_{m})$。


\subsection*{验证 (Verification)}

签名的验证按以下方式进行：

\begin{enumerate}
    \item 对所有 $j \in \{1,...,m\}$ 检查 $l \tilde{K}_j \stackrel{?}{=} 0$。\footnote{在 Monero 中，我们只为主 key image 检查 $l*\\tilde{K}_1 \\stackrel{?}{=} 0$。辅助密钥将存储为 $(1/8)*\\tilde{K}_j$，在验证时乘以8（回顾第 \ref{elliptic_curves_section} 节），这更高效。方法差异是一种实现选择，因为可链接的 key image 非常重要，不应被激进地修改，而另一种方法在之前的协议版本中使用过。}

    \item 为 \(i \in \{1, 2, ..., n\}\) 计算聚合公钥 $W_i$ 和聚合 key image $\tilde{W}$\vspace{.175cm}
    \begin{align*}
    W_i &= \sum^{m}_{j=1} \mathcal{H}_n(T_j, \mathcal{R}, \tilde{K}_1,...,\tilde{K}_{m})*K_{i,j}\\
    \tilde{W} &= \sum^{m}_{j=1} \mathcal{H}_n(T_j, \mathcal{R}, \tilde{K}_1,...,\tilde{K}_{m})*\tilde{K}_j
    \end{align*}{}

	\item 对于 \(i = 1, ..., n\) 计算，将 \(n + 1 \rightarrow 1\)，\vspace{.175cm}
	\[c_{i+1} = \mathcal{H}_n(T_c, \mathcal{R}, \mathfrak{m}, [r_i G + c_i W_i], [r_{i} \mathcal{H}_p(K_{i,1}) + c_i \tilde{W}])\]

	\item 如果 \(c_1 = c'_1\)，则签名有效。
\end{enumerate}


\subsection*{为什么有效 (Why it works)}

在这种简洁签名中最大的危险是密钥抵消，即报告的 key image 不是合法的，但仍然求和得到合法的聚合值。这就是聚合系数 $\mathcal{H}_n(T_j, \mathcal{R}, \tilde{K}_1,...,\tilde{K}_{m})$ 发挥作用的地方，它将每个密钥锁定到其预期值。我们将追踪伪造 key image 的循环影响作为练习留给读者（也许可以先想象这些系数不存在）。辅助 key image 是证明主 image 合法性的副产品，因为包含所有私钥的聚合私钥 $w_{\pi}$ 被应用于基底点 $\mathcal{H}_p(K_{\pi,1})$。


\subsection*{可链接性 (Linkability)}

如果私钥 \(k_{\pi, 1}\) 被重复使用来制作任何签名，签名中提供的对应主 key image \(\tilde{K}_1\) 将揭示这一点。辅助 key image 被忽略，因为它们的存在只是为了促进 CLSAG 的"简洁"部分。


\subsection*{空间需求 (Space requirements)}

我们存储 (1+$n$) 个整数，有 $m$ 个 key image，并使用 $m*n$ 个公钥。

%注意：在 Monero 中，CLSAG 密钥前缀略有不同。
%\[c_{i+1} = \mathcal{H}_n(T_c, one-time addresses, output commitments, pseudo output commitment, \mathfrak{m}, [r_i G + c_i W_i], [r_{i} \mathcal{H}_p(K_{i,1}) + c_i \tilde{W}])\]
